\begin{center}\begin{tikzpicture}[font=\small,>=stealth,line join=round]
\draw[gray!40,rounded corners=3pt] (0,0) rectangle (14.2,8.3);\node at (7.1,7.700000000000001) {例2：顶点在中点，两端同时达到最大值};
\begin{scope}[shift={(2.988888888888889,2.7115384615384612)},x=1.111111111111111cm,y=0.3538461538461538cm]\draw[->,gray] (-1.7,0)--(3.7,0) node[right]{$x$};\draw[->,gray] (0,-3)--(0,10) node[above]{$y$};\draw[blue,thick,very thick,domain=-1:3,samples=180] plot(\x,{2*(\x-1)^2-1});\fill[orange!85!black] (-1,7) circle(2.2pt);\fill[orange!85!black] (1,-1) circle(2.2pt);\fill[orange!85!black] (3,7) circle(2.2pt);\draw[violet,dashed] (1,-.8)--(1,9);\node[above] at (-1,7){$(-1,7)$};\node[above] at (3,7){$(3,7)$};\node[below] at (1,-1){$(1,-1)$};\end{scope}\node[anchor=west,align=left,text width=5.0cm] at (8.4,4.0) {$f=2(x-1)^2-1$\\[9pt]$x\in[-1,3]$\\[9pt]近处：$x=1$，最小$-1$\\[9pt]两端到轴距离都为2\\[9pt]最大值 $f(-1)=f(3)=7$};
\node[text width=12.6cm,align=center] at (7.1,.65) {同一最大值可以在两个不同位置取得；这里两个端点关于轴对称。};
\end{tikzpicture}\end{center}
